\documentclass{article}
\usepackage{graphicx}
\usepackage{amsmath}
\usepackage{amssymb}
\usepackage{amsfonts}
\usepackage{hyperref}
\usepackage{exercises}
\usepackage{indentfirst}
\usepackage{circledsteps}
\usepackage{amsthm}
\usepackage{csquotes}
\usepackage{yhmath}
\usepackage{comment}
\usepackage{placeins}
\usepackage{geometry}
\usepackage{hyperref}
\usepackage{tikz}
\newcommand*{\oldneg}{\mathord{\sim}}
\newcommand{\Z}{\mathbb{Z}}
\newcommand{\N}{\mathbb{N}}
\let\oldemptyset\emptyset

\theoremstyle{definition}
\newtheorem{definition}{Definition}
\newtheorem{theorem}{Theorem}[section]
\newtheorem{corollary}{Corollary}[theorem]
\newtheorem{lemma}[theorem]{Lemma}

\setlength{\parindent}{.5in}
\geometry{margin=0.5in}
\setlength{\parskip}{12pt}
\title{Chapter 11 Notes}
\begin{document}
\date{}
\maketitle

\section*{Section 0}

\subsection*{Notes}
Note that elementarily equivalence only requires pure formulas 
to be equivalent for the following reason: $A_1\ne A_2$ may be true, and if it is, 
then if $a\in A_1$ and $a\notin A_2$, then there are formulas either true or false in 
$(A_1,R_1)$ (formulas involving $a$ as a constant) where truth in $(A_2,R_2)$ isn't 
even defined.

\subsection*{Exercise 0.1}
Suppose $\pi$ is an isomoprhism from $(A_1,R_1)$ to $(A_2,R_2)$. 
We show by complete induction on the degrees of 
formulas that for any pure formula $\phi(x_1,...,x_n)$ 
and any elements $a_1,...,a_n$ of 
$A_1$, the sentence $\phi(a_1,...,a_n)$ 
is true in $(A_1,R_1)$ iff 
$\phi(\pi(a_1),...,\pi(a_n))$ is true in 
$(A_2,R_2)$, and, for any pure sentence $\phi$, 
$\phi$ is true in $(A_1,R_1)$ iff $\phi$ is true 
in $(A_2,R_2)$. 

In the base case, $Pab$ (for elements $a,b$ in $A_1$) is true in $(A_1,R_1)$ iff 
$aR_1b$ iff $\pi(a)R_2\pi(b)$ (because $\pi$ is an isomorphism) iff $P\pi(a)\pi(b)$ is true in 
$(A_2,R_2)$. Thus the base case holds.

Now suppose that the result holds for all formulas of degree less than $n$. That is, 
for any pure formula $\phi(x_1,...,x_m)$ (of a degree less than $n$) and any elements 
$a_1,...,a_m$ of $A_1$, $\phi(a_1,...,a_n)$ is true in $(A_1,R_1)$ iff
$\phi(\pi(a_1),...,\pi(a_m))$ is true in $(A_2,R_2)$, and for any pure sentence $\phi$, 
$\phi$ is true in $(A_1,R_1)$ iff $\phi$ is true in $(A_2,R_2)$. On the one hand, we consider a 
pure formula with free variables of degree $n$. That is, we consider some pure formula 
$\phi(x_1,...,x_m)$ and elements $a_1,...,a_m$ of $A_1$. There are then multiple cases.

In the first case, $\phi(x_1,...,x_m)=(\psi\wedge\gamma)$ for some pure formulas 
$\psi$ and $\gamma$ of degrees less than $n$. Either (1) 
neither $\psi$ nor $\gamma$ are sentences; (2) $\psi$ is a sentence and 
$\gamma$ is not; (3) $\psi$ is not a sentence and $\gamma$ is a sentence. We consider the first 
case - the other two are nearly identical. In that case, because neither $\psi$ nor 
$\gamma$ are sentences, $\psi=\psi(y_1,...,y_k)$ for some $k<n$.
Similarly $\gamma=\gamma(z_1,...,z_l)$ for some $l<n$. Then 
$\phi(x_1,...,x_m)=(\psi(y_1,...,y_k)\wedge\gamma(z_1,...,z_l))$. Then 
$\phi(a_1,...,a_m)=(\psi(b_1,...,b_k)\wedge\gamma(c_1,...,c_l))$ where 
each $b_i=a_j$ for some $j$ (and each $c_i=a_j$ for some $j$). Then 
$\phi(a_1,...,a_m)$ is true in $(A_1,R_1)$ iff 
$(\psi(b_1,...,b_k)\wedge\gamma(c_1,...,c_l))$ is true in $(A_1,R_1)$ iff 
$(\psi(\pi(b_1,...,\pi(b_k)))\wedge\gamma(\pi(c_1),...,\pi(c_l)))$ is true in 
$(A_2,R_2)$ (by the induction assumption) iff $\phi(\pi(a_1),...,\pi(a_m))$ is true 
in $(A_2,R_2)$.

The second case, where $\phi(x_1,...,x_m)=\oldneg\psi(x_1,...,x_m)$ 
(for some formula $\psi(x_1,...,x_m)$, is similar to the first case.

The third case is when $\phi(x_1,...,x_m)=(\exists x)\phi(x,x_1,...,x_m)$. Then 
$(\exists x)\phi(x,x_1,...,x_m)$ is true iff $\phi(a,a_1,...,a_m)$ is true for some 
$a\in A_1$, iff $\phi(\pi(a),\pi(a_1),...,\pi(a_m))$ is true in $(A_2,R_2)$ iff 
$(\exists x)\phi(x,\pi(a_1),...,\pi(a_m))$ is true in $(A_2,R_2)$.

On the other hand, we consider a pure sentence $\phi$ of degree $n$. The cases here are nearly 
identical to the above three cases - we give the case for existential quantification.

In that case, $\phi=(\exists x)\psi(x)$ for some $\psi$. Then
$\phi$ is true in $(A_1,R_1)$ iff $\psi(a)$ is true in $(A_1,R_1)$ for some 
$a\in A_1$ iff $\psi(\pi(a))$ is true in $(A_2,R_2)$ (by the induction assumption) iff 
$(\exists x)\psi(x)$ is true in $(A_2,R_2)$.

\end{document}

 
